\section{Effective division by \texorpdfstring{$1-\sigma$}{1-sigma}}
\label{sec:division}

The geometric argument supplies divisor classes that are divisible by
$\lambda=1-\sigma$ over the working field.  We now construct a divisor
representing such a division.  Its only nonlinear instruction is the
promised norm equation solved in Theorem~\ref{norm:solve}.  A
ramification point fixes the otherwise undetermined constant in that
equation.

\begin{samepage}
\begin{proposition}[Class division]\label{div:divide}
Let $k$, $q$, $F$, $C$, and $\sigma$ be as in
Proposition~\ref{div:arithmetic}, and suppose the auxiliary primary-group
data required by Theorem~\ref{norm:solve} are available.  Let $P=(x,0)$
be a specified $k$-rational ramification point.  Given a represented
degree-zero divisor $D$ satisfying
$\operatorname{char}k>(N+2\|D\|+1)^2$, under the promise that
\[
 [D]\in(1-\sigma)\operatorname{Pic}^0(C)(k),
\]
there is a deterministic procedure $\operatorname{LambdaDivide}(D,P)$
returning a degree-zero divisor $D'$ with
\[
 \|D'\|\leq2g,\qquad (1-\sigma)[D']=[D].
\]
It uses polynomially many field operations in the representation size,
$q,N,\|D\|$, and $\log|k|$, with the uniform norm bound from
Theorem~\ref{norm:solve}.  It is compatible with simultaneous
execution over the algebra of the unknown roots of $F$.
\end{proposition}
\end{samepage}

The procedure consists of three steps.
\begin{enumerate}
\item Compute $j=h^+/h^-$ by Lemma~\ref{div:compression}.  Put
$b=v_{X=x}(j)$ and normalize $G=cj$ so that $G/F^b$ has residue $1$
at $X=x$.
\item Use Theorem~\ref{norm:solve} to find $a$ with
$\operatorname{Norm}(a)=G$.  Set $E=D-\operatorname{div}(a)$, and form
the coefficientwise maximum $V$ of
$0,E,E+\sigma E,\ldots,E+\cdots+\sigma^{q-2}E$.
\item Return $\operatorname{ReduceDivisor}(V-(\deg V)P)$.
\end{enumerate}
We prove first that the norm equation in the second step is solvable,
and then that the maximum gives the required class division.

\begin{proof}
Lemma~\ref{div:compression} gives $j=h^+/h^-\in k(X)^*$ with
\begin{equation}\label{div:dividej}
 \operatorname{div}_C(j)=\sum_{a=0}^{q-1}\sigma^aD.
\end{equation}
Put $b=v_{X=x}(j)$.  Since $F$ has a simple root at $x$, the rational
function $j/F^b$ is a unit there, even when $b$ is negative.  Compute its
nonzero residue and set
\begin{equation}\label{div:normalization}
 c=\left(\left.\frac{j}{F^b}\right|_{X=x}\right)^{-1},
 \qquad G=cj.
\end{equation}
Thus $G/F^b$ has residue $1$ at $x$.

We first justify that $G$ is a norm from $\mathcal L=k(C)$ to $k(X)$,
which is the promise needed by the norm algorithm.  Choose a rational
divisor $D_0$ representing a class with $(1-\sigma)[D_0]=[D]$.
Lemma~\ref{geom:descent} supplies rational representatives and rational
functions for rational principal divisors.  There is therefore an
$a_0\in\mathcal L^*$ with
\[
 \operatorname{div}(a_0)=D-(1-\sigma)D_0.
\]
Taking orbit sums and using~\eqref{div:dividej} shows
\[
 \operatorname{Norm}_{\mathcal L/k(X)}(a_0)=c_0j
 \quad\text{for some }c_0\in k^*.
\]
At $P$ the cover is totally ramified, with residue field $k$, so the norm
valuation formula gives $v_P(a_0)=b$.  Also
\[
 \operatorname{Norm}_{\mathcal L/k(X)}(Y)=F
\]
because $q$ is odd.  The function $u=a_0/Y^b$ is thus a unit at $P$ and
has norm $c_0j/F^b$.  Every deck transformation acts trivially on its
residue, whence
\[
 \left.\frac{c_0j}{F^b}\right|_{X=x}=u(P)^q.
\]
By~\eqref{div:normalization}, the left side is $c_0/c$.  It follows that
$c/c_0$ is a $q$th power in $k^*$, and hence that $G=cj$ is a norm.
This argument is solely a proof of the promise; neither $D_0$ nor $a_0$
is an input to the algorithm.
Moreover, Lemma~\ref{div:compression} bounds the numerator and denominator
degrees of $G$ by $\|D\|$, so the stated characteristic hypothesis implies
the one in Theorem~\ref{norm:solve}.

Apply Theorem~\ref{norm:solve} to obtain $a\in\mathcal L^*$ with
$\operatorname{Norm}(a)=G$, and put
\[
 E=D-\operatorname{div}(a).
\]
Equation~\eqref{div:dividej} gives
$\sum_{a=0}^{q-1}\sigma^aE=0$.  Define the partial sums and their
coefficientwise maximum by
\begin{equation}\label{div:maximum}
 S_0=0,\qquad S_j=\sum_{a=0}^{j-1}\sigma^aE\ (1\leq j\leq q),
 \qquad V=\max_{0\leq j<q}S_j.
\end{equation}
Since $S_q=S_0=0$ and $\sigma S_j=S_{j+1}-E$, cyclically indexing the
same list shows
\[
 \sigma V=V-E,\qquad (1-\sigma)V=E.
\]
The maximum is effective because the list includes $0$.  As $P$ is
fixed by $\sigma$, the divisor
$V-(\deg V)P$ has degree zero and satisfies
\[
 (1-\sigma)\bigl[V-(\deg V)P\bigr]=[E]=[D].
\]
Return $\operatorname{ReduceDivisor}(V-(\deg V)P)$.  This proves the
asserted class identity and absolute-degree bound.

For completeness, every instruction in this description is effective in
the required representation.  The valuation and residue in
\eqref{div:normalization} are obtained by polynomial division at $X-x$.
The numerator and denominator of $j$ have degrees at most $\|D\|$;
multiplication by $c$ does not increase them.  Theorem~\ref{norm:solve}
therefore receives a rational function of that degree bound.  If its
output $a$ has coefficient and denominator degrees at most $E_a$, the
pole bound from Section~\ref{sec:divisors} gives
\[
 \|\operatorname{div}(a)\|\leq 4q(E_a+N).
\]
Thus $\|E\|\leq\|D\|+4q(E_a+N)$; the sum of the absolute degrees of
the partial sums in~\eqref{div:maximum} is at most $q^2\|E\|$.
The maximum and final reduction consequently have polynomial size.
Their ideal operations are precisely those of
Proposition~\ref{div:arithmetic}.  Scalar pivots, valuations, and all
other choices are determined by zero tests.  If an instruction differs
between root components, simultaneous execution returns the resulting
factor; otherwise the entire construction proceeds over the product
algebra.  This proves the complexity and simultaneous-execution claims.
\end{proof}
